\documentclass[12pt]{article}
\usepackage{hyperlatex}

%\usepackage{enumitem}
\usepackage{ulem}
\usepackage[compact,medium,sc]{titlesec}
\usepackage{booktabs}
\usepackage[margin=0.5in]{geometry}
\usepackage[pdftex]{graphicx}
\usepackage{multirow}
\usepackage{pslatex}
\usepackage{color}

% paragraph style

\setlength\parskip{10pt}
\setlength\parindent{0in}

\W\newcommand{\sf}{\bf}

\htmltitle{CS100: Project 1}
\htmldepth{1}
\htmladdress{\xlink{lusth@cs.ua.edu}{mailto: lusth@cs.ua.edu}}
\htmlcss{lusth.css}

\title{CS100: Project 1\\
\date{Revision Date: \today}}

\begin{document}

\maketitle

\thispagestyle{empty}

\W\subsubsection*{\xlink{Printable Version}{project1.pdf}}
\W\htmlrule

%Spring 2013

\section*{Preamble}

\begin{quote}
{\it A million here, a million there, and pretty soon you're talking real money -- Senator Everett Dirksen}
\end{quote}

%% Great quote!

One of the original purposes of the computer was to quickly perform
computations that took humans significant time and energy.
This program will do exactly that, calculating a result that could be
done with paper-and-pencil.  However, this program will allow us to
perform this calculation quickly and easily for a wide range of input values.

Imagine you are sitting in Bryant-Denny stadium and enduring a long television
timeout.  Your mind starts to wander.  Looking below you, you wonder how much
liquid (water) it would take to fill Bryant-Denny to your row.
You remember your high school science teacher making you figure out the volume of
an Olympic-sized swimming pool (660,000 gallons).  You wonder how the stadium compares.

Your task is to take some basic information about a stadium (think Bryant-Denny)
and do some calculations that will provide a useful answer to the question above.
That is, if you are on row {\bf N} of Bryant-Denny,
how many gallons (rounded to the nearest gallon)
of liquid would it take to fill the stadium to your feet?
We're in Alabama, so we're using standard U.S. gallons, not the imperial gallon or
the U.S. dry gallon.

\section*{Where to Start}

The first thing you should do is to setup a directory in which you will work.
Run the following commands in a terminal window, one at a time:

\begin{verbatim}
    cd
    cd cs100
    mkdir projects/project1
    cd projects/project1
\end{verbatim}

As a quick reminder, the first two commands move you into your {\it cs100}
directory, while the next two create and move you into a project directory.

After this, you need to ask yourself two questions.  First, what input do you
get when you run the program, the second is how you should model the stadium
to perform the calculations necessary to solve this problem.

\subsection*{What Input Do You Need?}

To answer the first question, your program will be given six pieces of information.

\begin{itemize}
\item Length - an integer value representing the length of the stadium floor (in yards).
This length includes the field, the endzones, plus any space beyond the endzones.
\item Width - an integer value representing the width of the stadium floor (in yards).
This width includes the field, the sidelines, plus any space beyond the sidelines.
\item Elevation - an integer value representing the height of the bottom row 
in the stadium (in feet).  Normally the bottom row starts several feet about the field
level for line-of-sight considerations.
\item Rise - an integer value representing the rise in height between rows in the
stadium (in inches).
\item Elevation angle -
an integer value that represents the angle (in degrees) of elevation for
the seats within the stadium.  We will assume the entire stadium has the same elevation
angle for all rows.  This is not normally the case.  If you want to learn more about
this, then do a web search for {\it C-value} and {\it sightline}.
\item Row - an integer representing the row where you are sitting in the stadium.
The very bottom row is considered {\bf row 0} to be consistent with counting in
{\tt C} programming.
\end{itemize}

\subsection*{How To Model the Stadium?}

To answer the second question,
you need to figure out how to compute the volume of space below you in the
stadium.  Here is a simplification of the stadium layout that
you must use for this project:

\begin{center}
\W\htmlimage[scale=0.50]{top-view.png}
\T\includegraphics{top-view.png}
\end{center}

The colors represent the various chunks that make up the total
volume.
You {\bf must} compute the volumes of these chunks
individually, via separate functions, and then add the
all return values up to get the total
volume.  Note: these functions should not print anything, just
return the volumes calculated. Chunks to consider include:

\begin{itemize}
\item the volume of the field itself (blue region),
extended vertically until it reaches the level of your row.  This is a
rectangular solid.
\item the volume of the two sidelines (green regions),
from the bottom row to your row.  This is a triangular wedge shape.
\item the volume of the two endzones (red regions),
from the bottom row to your row.  This is a triangular wedge shape.
\item the volume in the four corner areas (pink regions),
from the bottom row to your row.  The four pink regions can be
thought of as four pieces of a right circular cone (upside down).
\end{itemize}

\section*{Understanding How to Model the Stadium}

So how do you calculate the number of gallons of water that it would take to fill
Bryant-Denny to your feet?  First, we recommend you don't tackle the problem in terms
of gallons of water, but in terms of cubic feet.  Figure out how many cubic feet are
below your row.  Once you have this, it is easy to convert from cubic feet to gallons.

\subsection*{Simplifying the Model}

The volume below you is an irregular shape, but we 
will make some simplifying assumptions for our calculations.  These assumptions are:

\begin{itemize}
\item The field is rectangular, with a known length and width
\item The elevation from the field to the bottom row of the stands is a vertical rise
\item The stands start at the same height everywhere (sidelines, endzones, corners)
\item The stands on both sidelines have the same characteristics and dimensions
\item The two endzone stands have the same characteristics and dimensions
\item The corners (between the sideline stands and endzone stands) all have the same 
characteristics and dimensions
\end{itemize}

\subsection*{Modeling the Different Parts}

When looking at the total number of cubic feet, don't try and solve this all at once.
Instead, tackle it in parts.  Then simply add the various parts to get a total
volume in cubic feet.  Specifically:

\begin{itemize}
\item Model the field itself as a rectangular solid, you just need to know the length
and width and height to calculate its volume.
\item Model both the sideline stands and the endzone stands as a wedge.  That is,
a right triangle that has a known thickness. The thickness of the wedge corresponds to the
length of the sideline stands.  To construct the right triangle, you'll need to remember
some basic trig functions.  You have a right triangle, know the angle of elevation,
and you know the rise in height.  You've got two angles and a side, you just need
to find the length of the other side for your volume calculations.
\item Model the endzone stands the exact same way you did the sideline stands.  The
height and angle of elevation is the same, the only difference is the thickness of this
wedge.  It is now the width of the field instead of the length of the field.
\item Model each of the four corners as one-fourth of a right circular cone.  In our case,
the cone is upside down from what you see in most pictures.  There are some
basic formulas to calculate the volume of a right circular cone that you'll need to
find for this program.  You need to figure out the height and radius for this
right circular cone.  The height is the same as the sidelines and endzone, and the
angle of elevation plus rise between rows should allow you to calculate the radius.
\end{itemize}

One final comment on this.  If you don't remember what these basic shapes look like, 
search online for images of rectangular solid and right circular cone and triangle wedge.
Just remember to think of the right circular cone upside down.

As mentioned previously, we {\it strongly} recommend that you structure your
program so that your calculations are done in pieces.  As you compute the
various pieces of information that you need, print them out to confirm their
accuracy.

\section*{Programming This Task}

Do {\bf not} approach this problem as one long main program.  If you do, you
will lose {\bf significant} points.  Instead, break the task into
parts and have a separate function for each part.  At a minimum, we suggest:

\begin{itemize}
\item A function to read the six input values
\item A function to print out the six input values that you received
\item A function to calculate the volume of the stadium below a given row.
Note that you must use several {\it helper} functions in
this function.  These helper functions would:
\begin{itemize}
\item Calculate the volume above the field
\item Calculate the volume in the sideline (and endzone) stands
\item Calculate the volume of the four corners
\end{itemize}
\end{itemize} 

In addition to these functions, there are several {\it helper}
functions that can be used at various points in your program.  These
would include:

\begin{itemize}
\item A function to convert from yards to feet
\item A function to convert from inches to feet
\item A function to convert from cubic feet to gallons
\item A function to convert from degrees to radians (since {\tt C} uses 
radians in its various mathematical functions -- {\it sin, cos, tan, etc.})
\end{itemize}

When programming, do not try and write in one sitting.  Instead, we recommend
that you build this in steps.  Try getting things to work in the order listed below:

\begin{enumerate}
\item Get a program that reads in the six values and print them out
\item Get a program that reads in six values, prints them out, and then
calculates the rectangular volume of the field (to the given row).  Ignore the volume
in the stands themselves.
\item Get a program that reads in six values, prints them out,
and then calculates
the rectangular volume of the field and the volume of the triangular wedges on the
sidelines.  Ignore the volume in the four corners.
\item Get a program that reads in six values, prints them out, 
nd then calculates
the total volume of the field (to that given row).
\item Get a program that reads in six values, prints them out, 
prints out the intermediate and final volumes in
cubic feet,
and finally prints out the final volume to the nearest gallon.
\end{enumerate}

\section*{Compliance Instructions}

Your program should be named {\tt stadium.c}.

Compile and test your program, naming the executable {\it stadium}, by entering
the following commands and input:

\begin{verbatim}
    $ gcc -Wall -g stadium.c -o stadium -lm
    $ stadium
    length? 0
    width? 0
    elevation? 0
    rise? 0
    angle? 45
    row? 0
    --------------------------
    length is 0     width is 0       elevation is 0
    rise is 0       angle is 45     row is 0
    --------------------------
    field volume is 0.000000
    sideline volume is 0.000000
    endzone volume is 0.000000
    corner volume is 0.000000
    total volume in cubic feet is 0.000000
    ---------------------------
    0 gallons
\end{verbatim}

Your program should print out clearly the intermediate results
and the final result. Add a little more explanatory text for the
output than shown above (e.g. units).

You must generate at least one {\tt test.dat} file that contains a set 
of {\it reasonable} data for this project.   Test your data with the command:

\begin{verbatim}
    cat test.dat | stadium
\end{verbatim}

\color{red}
If your code fails with a runtime error while running this test
or if you fail to provide a working {\it test.dat},
then you will receive a zero for this assignment.
\color{black}

Note that your answers do not have to be correct for your program to be graded,
only that the program does not crash. Of course, correct answers will yield 
a much higher grade.

\section*{Checklist}

Did you:

\begin{itemize}
\item
    name your program {\it stadium.c}?
\item
    use consistent units for calculating your volumes?
\item
    use lots of functions?
\item
    have not functions printing other than  your main function and your procedure for printing the input values?
\item
    print out the intermediate results?
\item
    round to the nearest gallon?
\item
    create your {\it test.dat} file and test your program with it?
\item
    check your {\it test.dat} results on the forum?
\end{itemize}

Note: this list is not exhaustive! There may be other specifications
and requirements not listed here.

\section*{Submission Instructions}

Change to the {\it project1} directory containing your assignment.  Do an
{\it ls} command. You should see something like this:

\begin{verbatim}
    stadium  stadium.c  test.dat
\end{verbatim}

Extra files are OK, as long as you are in the correct directory. 
Submissions from the wrong directory will be penalized.

Submit your program like this:

\begin{verbatim}
    submit cs100 YYY project1
\end{verbatim}

Remember to replace \verb!YYY! with your instructor name
(brown,cordes,lusth).

\section*{Due Date}

The due date for this assignment can be found on the class
\xlink{schedule}{../schedule.html}.

\end{document}
